% !TEX root = ../main.tex
\chapter{AURORA: projeto, processadores e árvore de arquivos}
\label{cap:07-aurora-projeto-arvore}

Este capítulo descreve o modelo de projeto da \tool{AURORA} — como um projeto
é representado no disco e em memória — e a navegação pela árvore de arquivos.
A fonte da verdade é o código-fonte do renderer (\path{js/project/},
\path{js/tree/}, \path{js/processors/}) e dos handlers do processo principal
(\path{main/ipc/project.js}, \path{main/ipc/files.js}, \path{main/recents.js}).
O fluxo de compilação que consome o que está aqui descrito (botões Verilog,
Wave, PRISM, \cmm) é assunto do \autoref{cap:08-aurora-toolchain} e não é
re-documentado aqui.

\section{Visão geral do modelo de projeto}

Um projeto da \tool{AURORA} é uma pasta no disco que contém um arquivo
\path{.spf} (\enquote{SAPHO Project File}) de mesmo nome que a pasta. O
\path{.spf} consolida toda a configuração canônica do projeto: metadados de
ciclo de vida (nome, datas, máquina, versão do app) e a \emph{estrutura}
(listas de arquivos, top-level, testbench-topo, processadores). Em torno desse
arquivo orbitam subpastas — uma por processador \tool{SAPHO} — e os arquivos
\file{.v}/\file{.sv} importados pelo usuário.

A representação em memória é dividida entre o processo principal (que faz o IO
no disco) e o renderer (que mantém o estado da UI). Três peças de software
governam o estado:

\begin{description}[leftmargin=2.2em,style=nextline]
  \item[\file{ProjectStore} (\path{js/project/project_store.js})] dono único de
    \enquote{qual projeto está aberto} no renderer — guarda o
    \lstinline|projectPath| (pasta) e o \lstinline|spfPath| (arquivo \path{.spf}).
  \item[\file{SpfStore} (\path{js/project/spf_store.ts})] escritor único e
    serializado do \path{.spf} no lado do renderer; faz leitura pura e escrita
    atômica do bloco \lstinline|structure|.
  \item[\file{state.currentOpenProjectPath} (\path{main/state.js})] no processo
    principal, é a única fonte de verdade de \enquote{qual \path{.spf} está
    aberto}; a pasta do projeto é sempre derivada por \lstinline|path.dirname|.
\end{description}

\begin{nota}
A pasta do projeto NUNCA é armazenada em paralelo ao \path{.spf} no processo
principal: ela é sempre derivada de \lstinline|path.dirname(spfPath)| sob
demanda. O comentário \enquote{A4} em \path{main/ipc/project.js} registra essa
decisão — evita o problema clássico de dois globais que divergem.
\end{nota}

\section{O arquivo de projeto \texttt{.spf}}
\label{sec:spf-formato}

O \path{.spf} é um documento JSON com dois blocos de topo: \lstinline|metadata|
e \lstinline|structure|. O esqueleto é produzido pela classe
\lstinline|ProjectFile| em \path{main/ipc/project.js} no momento da criação.

\subsection{Estrutura completa}

A tabela seguinte lista todos os campos definidos no código. As interfaces
\lstinline|SpfDocument|, \lstinline|SpfStructure| e \lstinline|SpfMetadata| em
\path{spf_store.ts} são a referência do lado do renderer; o construtor de
\lstinline|ProjectFile| é a referência do lado do processo principal.

\begin{longtable}{p{0.30\linewidth} p{0.16\linewidth} p{0.46\linewidth}}
  \toprule
  \textbf{Campo} & \textbf{Bloco} & \textbf{Significado} \\
  \midrule
  \endhead
  \lstinline|projectName|   & metadata  & Nome do projeto (basename da pasta). Exibido na barra de título como \lstinline|<nome>.spf|. \\
  \lstinline|createdAt|     & metadata  & Timestamp ISO de criação. \\
  \lstinline|lastModified|  & metadata  & Timestamp ISO carimbado a cada escrita do \path{.spf}. \\
  \lstinline|lastOpened|    & metadata  & Timestamp ISO da última abertura (gravado em \lstinline|project:open|). \\
  \lstinline|computerName|  & metadata  & \lstinline|COMPUTERNAME| ou \lstinline|os.hostname()|. \\
  \lstinline|appVersion|    & metadata  & Versão do app (\lstinline|app.getVersion()|). \\
  \lstinline|projectPath|   & metadata  & Caminho absoluto da pasta do projeto. \\
  \midrule
  \lstinline|basePath|          & structure & Caminho absoluto da pasta do projeto; reconciliado na abertura. \\
  \lstinline|folders|           & structure & Subpastas registradas pela UI (criadas via \enquote{New Folder}). \\
  \lstinline|processors|        & structure & Lista de processadores \tool{SAPHO}; cada entrada é \lstinline|{ name }|. \\
  \lstinline|topLevelFile|      & structure & Top do design (modo \emph{check} do iverilog, flag \lstinline|-s|). \\
  \lstinline|testbenchFile|     & structure & Testbench-topo (modo \emph{simulação}, flag \lstinline|-s| do botão Wave). \\
  \lstinline|synthesizableFiles|& structure & Arquivos \file{.v} sintetizáveis; cada entrada tem \lstinline|path| (e \lstinline|name|). \\
  \lstinline|testbenchFiles|    & structure & Arquivos \file{.v} de testbench; mesma forma de entrada. \\
  \bottomrule
\end{longtable}

As constantes de \lstinline|STRUCTURE_DEFAULTS| em \path{spf_store.ts}
congelam (\lstinline|Object.freeze|) os valores-padrão de \lstinline|structure|:
\lstinline|basePath: ''|, \lstinline|folders: []|, \lstinline|processors: []|,
\lstinline|topLevelFile: ''|, \lstinline|testbenchFile: ''|,
\lstinline|synthesizableFiles: []| e \lstinline|testbenchFiles: []|. Na
leitura, os defaults são aplicados primeiro e os valores do disco depois
(\code{\{ ...STRUCTURE\_DEFAULTS, ...parsed.structure \}}), de modo que chaves
desconhecidas sobrevivem ao \emph{round-trip} (\enquote{unknown keys survive}).

\subsection{Exemplo de um \texttt{.spf}}

\begin{lstlisting}
{
  "metadata": {
    "projectName": "demo",
    "createdAt": "2026-06-21T12:00:00.000Z",
    "lastModified": "2026-06-21T12:34:56.000Z",
    "computerName": "NIPS-CERN-PC",
    "appVersion": "1.0.0",
    "projectPath": "C:\\Users\\dev\\Documents\\demo"
  },
  "structure": {
    "basePath": "C:\\Users\\dev\\Documents\\demo",
    "folders": [],
    "processors": [ { "name": "cpu" } ],
    "topLevelFile": "cpu\\Hardware\\cpu.v",
    "testbenchFile": "cpu\\Simulation\\cpu_tb.v",
    "synthesizableFiles": [
      { "name": "cpu.v", "path": "cpu\\Hardware\\cpu.v" }
    ],
    "testbenchFiles": [
      { "name": "cpu_tb.v", "path": "cpu\\Simulation\\cpu_tb.v" }
    ]
  }
}
\end{lstlisting}

\subsection{Normalização de caminhos: relativo no disco, absoluto em memória}
\label{sec:spf-paths}

O \file{SpfStore} grava os caminhos de arquivo \emph{relativos} ao
\lstinline|basePath| quando o arquivo está dentro da pasta do projeto, e
\emph{absolutos} quando está fora. Tudo que é exposto fora do \file{SpfStore}
(via \lstinline|read|/\lstinline|update|) é \emph{sempre} absoluto — o chamador
não precisa conhecer essa normalização. As funções relevantes em
\path{spf_store.ts} são:

\begin{itemize}[leftmargin=1.6em]
  \item \lstinline|isAbsolutePath| — reconhece caminhos Windows
    (\lstinline|C:\|, \lstinline|\\server\|) e POSIX (\lstinline|/|).
  \item \lstinline|toRelativeIfInside| — converte para relativo se o absoluto
    estiver sob o \lstinline|baseDir| (comparação \emph{case-insensitive},
    idioma Windows); caso contrário retorna o absoluto intacto.
  \item \lstinline|toAbsoluteFromBase| — junta um relativo com o
    \lstinline|baseDir| na leitura.
  \item \lstinline|expandStructurePaths| / \lstinline|contractStructurePaths| —
    aplicam a conversão sobre os campos de caminho (\lstinline|topLevelFile|,
    \lstinline|testbenchFile|, \lstinline|synthesizableFiles[].path|,
    \lstinline|testbenchFiles[].path|).
\end{itemize}

\begin{nota}
\path{.spf} antigos com caminhos absolutos continuam sendo lidos
(\lstinline|isAbsolutePath| faz \emph{bypass} do \lstinline|resolve|); o
primeiro \lstinline|save| migra automaticamente para a forma relativa. Isso
permite copiar/mover a pasta do projeto sem quebrar as referências internas.
\end{nota}

\subsection{Parsing tolerante}

A leitura do \path{.spf} usa \lstinline|parseSpfTolerant| (presente tanto em
\path{spf_store.ts} quanto em \path{main/ipc/project.js}): tenta
\lstinline|JSON.parse| estrito primeiro e, em caso de falha, faz uma única
passagem leniente que remove BOM, comentários (\lstinline|//| e
\lstinline|/* */| fora de strings) e vírgulas finais (\emph{trailing commas}).
Um arquivo recuperável não é silenciosamente resetado para os defaults; só um
arquivo genuinamente ilegível cai no fallback (defaults), com aviso no console.

\section{O \texttt{SpfStore}: escritor único e serializado}
\label{sec:spfstore}

O \file{SpfStore} (\path{js/project/spf_store.ts}, compilado por
\lstinline|tsc| em \path{spf_store.js}) é o único ponto de escrita do
\path{.spf} no renderer. Sua API espelha um antigo \lstinline|ProjectConfigStore|:

\begin{description}[leftmargin=2em,style=nextline]
  \item[\lstinline|read(spfPath)|] devolve só o bloco \lstinline|structure|
    (com defaults aplicados, caminhos expandidos para absolutos).
  \item[\lstinline|readFull(spfPath)|] devolve \lstinline|metadata| +
    \lstinline|structure|.
  \item[\lstinline|update(spfPath, mutator)|] leitura-mutação-escrita atômica;
    o \lstinline|mutator| recebe o \lstinline|structure| e o \lstinline|metadata|
    é preservado com \lstinline|lastModified| recarimbado.
\end{description}

\subsection{Serialização das escritas}

Cada \path{.spf} tem uma cadeia de promises própria
(\lstinline|writeChainByPath: Map|): chamadas de \lstinline|update| para o mesmo
caminho serializam em ordem de chegada (\lstinline|prev.then(...)|), de modo que
duas mutações no mesmo \emph{tick} não operem sobre o mesmo \emph{snapshot} pré-escrita.

\subsection{Coalescência de leituras}

Para evitar \emph{fanout} de \lstinline|readFile|/\lstinline|JSON.parse|, o
\file{SpfStore} coalesce leituras em voo (\lstinline|inFlightReads: Map|): se
vários consumidores chamam \lstinline|read| no mesmo tick, todos compartilham a
mesma promise. O mapa é limpo quando a promise é resolvida — não há cache
atravessando ticks (sem \emph{mtime} obsoleto). O \lstinline|update|, por outro
lado, sempre lê do disco com \lstinline|readRawUncached| (o \lstinline|writeChain|
já serializa, então pular o coalescing aqui não introduz \emph{race}).

\subsection{Evento \texttt{aurora:spf-changed}}

Após cada escrita, o \file{SpfStore} compara o \lstinline|structure| antes e
depois do \lstinline|mutator| (via \lstinline|JSON.stringify|). Somente quando o
\lstinline|structure| efetivamente mudou é disparado o evento de janela
\lstinline|aurora:spf-changed| (com \lstinline|detail: { spfPath, source }|).
Mutações \emph{no-op} não geram refresh — isso elimina a classe de laços
\enquote{save $\rightarrow$ event $\rightarrow$ refresh $\rightarrow$ load
$\rightarrow$ save (no-op) $\rightarrow$ \ldots}. Os assinantes desse evento
incluem a barra de status, o painel de configuração do processador e o
\file{file\_mode}.

\section{O \texttt{ProjectStore}: dono do projeto aberto}

O \file{ProjectStore} (\path{js/project/project_store.js}) mantém em variáveis
de módulo o \lstinline|projectPath| e o \lstinline|spfPath| do projeto aberto.
A API é \lstinline|setProject(spfPath, projectPath)|, \lstinline|clearProject()|,
\lstinline|getProjectPath()|, \lstinline|getSpfPath()|, \lstinline|hasProject()|
e \lstinline|subscribe(fn)|.

Por compatibilidade com dezenas de pontos de leitura legados, o store
\emph{espelha} seus valores para \lstinline|window.currentProjectPath| e
\lstinline|window.currentSpfPath| a cada mudança (\lstinline|mirrorToWindow|).
Código novo deve preferir os getters. Assinantes (\lstinline|subscribe|) são
notificados com um \emph{snapshot} \lstinline|{ projectPath, spfPath }| — o
controlador de views da árvore, por exemplo, usa isso para habilitar/desabilitar
a view \enquote{standard} conforme um projeto abre ou fecha.

\section{Ciclo de vida: criar, abrir e fechar projeto}

\subsection{Criar projeto}

O fluxo começa na UI: \path{js/ui/new_project_modal.js} valida o nome
(\lstinline|/^[a-zA-Z0-9_-]+$/|) e o caminho, monta
\lstinline|projectPath = <local>\<nome>| e
\lstinline|spfPath = <projectPath>\<nome>.spf|, e chama
\lstinline|electronAPI.createProjectStructure|. No processo principal, o handler
\lstinline|project:createStructure| (\path{main/ipc/project.js}):

\begin{enumerate}[leftmargin=1.8em]
  \item cria a pasta do projeto (\lstinline|fse.mkdir| recursivo);
  \item instancia \lstinline|new ProjectFile(projectPath)| e grava o
    \path{.spf} com \lstinline|metadata| + \lstinline|structure| iniciais;
  \item confirma a existência da pasta e do \path{.spf};
  \item lista o conteúdo da pasta para devolver ao renderer;
  \item registra o novo \path{.spf} nos recentes e reconstrói a \emph{jumplist}
    do Windows.
\end{enumerate}

A localização escolhida é lembrada em
\lstinline|aurora-last-new-project-location| (localStorage) para que a próxima
caixa de diálogo \enquote{New Project} abra na mesma pasta-pai. Em seguida o
modal chama \lstinline|projectManager.loadProject(spfPath)|.

\subsection{Abrir projeto}

Há vários gatilhos de abertura, todos convergindo para
\lstinline|loadProject| em \path{js/project/project_manager.js}:

\begin{itemize}[leftmargin=1.6em]
  \item botões \enquote{Open Project} (toolbar e tela de boas-vindas), que
    chamam \lstinline|electronAPI.showOpenDialog|;
  \item clique numa linha de projeto recente (tela de boas-vindas);
  \item abertura via \enquote{File > Open} / atalhos
    (\lstinline|onSimulateOpenProject|).
\end{itemize}

O handler do processo principal \lstinline|project:open|
(\path{main/ipc/project.js}) faz o trabalho pesado:

\begin{enumerate}[leftmargin=1.8em]
  \item Corrige o caminho se o \path{.spf} não existir no formato esperado
    (tenta \lstinline|<root>/<nome>/<nome>.spf|).
  \item Define \lstinline|state.currentOpenProjectPath = spfPath| (fonte de
    verdade no main).
  \item Registra nos recentes (\path{main/recents.js}) e reconstrói a
    \emph{jumplist}; chama também \lstinline|app.addRecentDocument|.
  \item Lê e faz parse tolerante do \path{.spf}, carimba
    \lstinline|metadata.lastOpened|.
  \item \textbf{Reconcilia o \lstinline|basePath|}: alinha sempre com
    \lstinline|path.dirname(spfPath)|. Se diferir, faz \lstinline|deepRemapPaths|
    de todo caminho absoluto que ainda apontava para a raiz antiga (caso de
    projeto copiado/movido) — ver \autoref{sec:rename}.
  \item Para cada processador, anexa \lstinline|exists| (existência da pasta no
    disco).
  \item Reescreve o \path{.spf} (já reconciliado) e devolve
    \lstinline|{ projectData, files, spfPath }|, enviando também os eventos IPC
    \lstinline|project:processorHubState| e \lstinline|project:processors| para
    a janela que originou a requisição.
\end{enumerate}

De volta no renderer, \lstinline|loadProject| orquestra a UI:

\begin{enumerate}[leftmargin=1.8em]
  \item resolve o \lstinline|basePath| de forma tolerante (várias formas de
    payload) e chama \lstinline|ProjectStore.setProject(spfPath, basePath)|;
  \item \lstinline|projectTreeManager.reset()| — \emph{clean slate} antes da
    nova árvore (uma troca direta projeto$\rightarrow$projeto não passa pelo
    fluxo de fechar);
  \item \lstinline|setAvailableProcessors(structure.processors)| — semeia a
    lista global de processadores ANTES de renderizar;
  \item atualiza o nome na UI, fecha todas as abas
    (\lstinline|TabManager.closeAllTabs|);
  \item \lstinline|projectTreeManager.activateTree()| popula e renderiza a
    árvore;
  \item inicia o \emph{watcher} de arquivos, registra o projeto nos recentes da
    tela de boas-vindas e habilita os botões de compilação;
  \item sintetiza \lstinline|aurora:spf-changed| (\lstinline|source: 'project-loaded'|)
    para forçar a reaplicação do estado derivado do \path{.spf} (testbench,
    top-level, processadores) em todos os assinantes;
  \item se o \path{.spf} referencia arquivos que sumiram do disco, gera o
    relatório \path{.aurora-missing-files.log} (aberto no Monaco como preview) e
    uma notificação; caso contrário, remove um log obsoleto de corrida anterior.
\end{enumerate}

\subsection{Fechar projeto}

\path{js/project/close_project.js} trata o botão de fechar: pede confirmação
(\lstinline|showDialog|), chama \lstinline|electronAPI.closeProject| (handler
\lstinline|project:close| no main, que zera \lstinline|currentOpenProjectPath| e
emite \lstinline|project:processors|/\lstinline|project:fileTree|/\lstinline|project:closed|
vazios), fecha todas as abas, limpa a interface
(\lstinline|clearProjectInterface| $\rightarrow$ \lstinline|treeView.clearAll()|
+ card de \enquote{no project}), chama \lstinline|ProjectStore.clearProject()|,
esquece o último projeto aberto (\lstinline|appInitializer.clearLastProject|) e
faz \lstinline|projectTreeManager.reset()| para que a próxima abertura dispare
um \emph{re-activate} completo.

\section{Processadores}
\label{sec:processadores}

Um processador \tool{SAPHO} corresponde a uma subpasta da raiz do projeto, com
três subpastas reconhecidas: \path{Software/}, \path{Hardware/} e
\path{Simulation/}. A lista canônica de processadores do projeto é
\lstinline|structure.processors| no \path{.spf} (cada entrada é
\lstinline|{ name }|); em memória, no renderer, a lista de nomes vive em
\lstinline|window.availableProcessors|.

\subsection{A lista global de processadores}

O módulo \file{processor\_list} (\path{js/project/processor_list.ts}, compilado
em \path{.js}) é o dono único de \lstinline|window.availableProcessors| e expõe:

\begin{itemize}[leftmargin=1.6em]
  \item \lstinline|setAvailableProcessors(list)| — substitui a lista inteira;
    aceita arrays de strings ou de \lstinline|{name}| e faz \emph{dedup}
    \emph{case-insensitive} (pastas no Windows são \emph{case-insensitive});
  \item \lstinline|getAvailableProcessors()| — sempre retorna array;
  \item \lstinline|addAvailableProcessor(name)| — acrescenta um nome se ainda
    não presente (match \emph{case-insensitive}).
\end{itemize}

\subsection{Criar processador}

A criação é feita pelo \emph{Processor Hub} (\path{js/processors/processor_hub.js}):
um formulário modal com validação ao vivo. As regras de validação são:

\begin{longtable}{p{0.30\linewidth} p{0.62\linewidth}}
  \toprule
  \textbf{Campo} & \textbf{Regra} \\
  \midrule
  \endhead
  \lstinline|name|         & \lstinline|/^[a-zA-Z0-9_-]+$/| (letras, números, hífen, sublinhado). \\
  \lstinline|nBits|        & Inteiro positivo; deve satisfazer a consistência abaixo. \\
  \lstinline|nbMantissa|   & Inteiro positivo. \\
  \lstinline|nbExponent|   & Inteiro positivo. \\
  \lstinline|gain|         & Potência de 2 (\lstinline|(num & (num-1)) === 0|). \\
  \lstinline|instructionStackSize| & Inteiro positivo. \\
  \lstinline|dataStackSize|        & Inteiro positivo. \\
  \lstinline|inputPorts|   & Inteiro positivo. \\
  \lstinline|outputPorts|  & Inteiro positivo. \\
  \bottomrule
\end{longtable}

A consistência de bits exige \lstinline|nBits === nbMantissa + nbExponent + 1|
(\lstinline|checkBitConsistency|). O botão \enquote{Generate} só fica habilitado
quando todos os campos passam (\lstinline|validateAll|).

No envio, o renderer chama \lstinline|electronAPI.createProcessorProject|; o
handler \lstinline|create-processor-project| (\path{main/ipc/project.js}):

\begin{enumerate}[leftmargin=1.8em]
  \item re-valida o nome (defesa contra \emph{path traversal}: o nome vai direto
    para \lstinline|path.join|, então é confinado a um único segmento);
  \item cria \path{<proc>/Software}, \path{<proc>/Hardware} e
    \path{<proc>/Simulation};
  \item escreve \path{Software/<proc>.cmm} com um cabeçalho de diretivas
    \tool{SAPHO} preenchido a partir do formulário;
  \item acrescenta \lstinline|{ name }| em \lstinline|structure.processors| do
    \path{.spf} (com \emph{dedup case-insensitive}) e reescreve o arquivo;
  \item emite o evento IPC \lstinline|processor:created| para o renderer.
\end{enumerate}

O cabeçalho \file{.cmm} gerado tem a forma:

\begin{lstlisting}[style=cmm]
#PRNAME <nome>
#NUBITS <nBits>
#NDSTAC <dataStackSize>
#SDEPTH <instructionStackSize>
#NUIOIN <inputPorts>
#NUIOOU <outputPorts>
#NBMANT <nbMantissa>
#NBEXPO <nbExponent>
#NUGAIN <gain>

void main()
{
    // ...
}
\end{lstlisting}

\begin{nota}
As diretivas \lstinline|#NUBITS|, \lstinline|#NBMANT| etc. no \file{.cmm} são a
fonte \emph{canônica} dos parâmetros do processador para o pipeline de
compilação. O \path{.spf} guarda apenas \lstinline|{ name }|; os números vivem
no \file{.cmm}. A semântica das diretivas é detalhada nos capítulos do
\tool{YANC}.
\end{nota}

\subsection{Renomear processador}
\label{sec:rename}

O handler \lstinline|rename-processor| (\path{main/ipc/project.js}) renomeia o
processador em todas as superfícies internas do \tool{SAPHO}, de forma atômica:

\begin{enumerate}[leftmargin=1.8em]
  \item libera os \emph{watchers} sob a raiz do projeto
    (\lstinline|releaseWatchersUnder|) — sem isso o \tool{chokidar} mantém um
    \emph{handle} e a renomeação da pasta falharia com EPERM/EBUSY no Windows;
  \item move a pasta \path{<root>/<old>} $\rightarrow$ \path{<root>/<new>}
    (\lstinline|moveWithRetry|; renomeação só de \emph{case} usa um salto por
    pasta temporária);
  \item renomeia os artefatos nomeados pelo processador:
    \path{Software/<old>.cmm}, \path{Software/<old>.asm},
    \path{Hardware/<old>.v}, \path{Simulation/<old>_tb.v};
  \item corrige a diretiva \lstinline|#PRNAME| no \file{.cmm} — \emph{apenas} a
    linha da diretiva; comentários e código do usuário não são tocados;
  \item atualiza a entrada em \lstinline|processors[]| preservando a
    configuração por processador;
  \item faz \lstinline|remapProcessorPath| de qualquer referência do \path{.spf}
    (\lstinline|topLevelFile|, \lstinline|testbenchFile|,
    \lstinline|synthesizableFiles|, \lstinline|testbenchFiles|) que apontava para
    dentro da pasta antiga;
  \item emite \lstinline|project:processors| e \lstinline|processor:renamed|.
\end{enumerate}

\subsection{Deletar processador}

O handler \lstinline|delete-processor| remove a pasta do processador do disco
(\lstinline|fse.remove|), filtra a entrada de \lstinline|structure.processors| no
\path{.spf} e emite \lstinline|project:processors| atualizado. No renderer, a
opção \enquote{Delete processor} aparece no menu de contexto do separador de
processador na árvore (\lstinline|_deleteProcessorByName| em
\path{project_tree_actions.js}, que chama \lstinline|electronAPI.deleteProcessor|).

\subsection{Agrupamento e identificação na árvore}

Na view de arquivos (\enquote{verilog picker}), os arquivos são agrupados por
processador. A função \lstinline|_getProcessorForFile| em \path{file_mode.js}
decide o \enquote{dono} de um arquivo: o caminho precisa cair em
\path{<projeto>/<proc>/{Hardware|Software|Simulation}/<arquivo>} e
\lstinline|<proc>| precisa estar em \lstinline|window.availableProcessors|
(comparação \emph{case-insensitive}). Cada grupo é precedido por um separador
horizontal rotulado com o nome do processador
(\lstinline|_createProcessorSeparator| em \path{project_tree_render.js}); o
grupo de arquivos importados que não pertencem a nenhum processador aparece sob
um separador \enquote{Imported} (apenas quando há também grupos de processador).

\section{As três views da árvore de arquivos}
\label{sec:views}

A árvore de arquivos da \tool{AURORA} pode se apresentar de três formas, todas
dentro do mesmo elemento \lstinline|#file-tree|:

\begin{description}[leftmargin=2em,style=nextline]
  \item[\enquote{verilog}] o \emph{verilog picker} — a view de arquivos do
    projeto, agrupada por processador, com indicadores de synth/testbench. É a
    view padrão e também hospeda os cartões de estado vazio.
  \item[\enquote{hierarchy}] a árvore de instâncias de módulo produzida pelo
    \tool{Yosys} (resultado de compilação).
  \item[\enquote{standard}] uma árvore simples de pastas/arquivos enraizada na
    pasta do projeto, lendo filhos sob demanda (\emph{lazy}).
\end{description}

\subsection{Subcontêineres físicos e o \texttt{tree\_view}}

O controlador de baixo nível \file{tree\_view} (\path{js/tree/tree_view.js})
cria três subcontêineres DOM fisicamente separados dentro de
\lstinline|#file-tree| — um por view — e governa a visibilidade por um único
atributo \lstinline|data-active-view| (a CSS em \path{css/tree/file_tree.css}
mostra/oculta com base nele). Cada \emph{renderer} escreve apenas no \emph{seu}
subcontêiner — eles não podem colidir.

\begin{lstlisting}
<div id="file-tree" data-active-view="verilog">
  <div class="tree-view tree-view-verilog">...</div>
  <div class="tree-view tree-view-hierarchy">...</div>
  <div class="tree-view tree-view-standard">...</div>
</div>
\end{lstlisting}

A API é \lstinline|getContainer(name)| (devolve o subcontêiner), \lstinline|setActive(name)|
(troca a view visível — uma única mudança de atributo, sem mutação de DOM),
\lstinline|getActive()|, \lstinline|clear(name)| e \lstinline|clearAll()|. O
array \lstinline|VIEW_NAMES| congela os três nomes.

\subsection{O controlador único de views}

O \file{file\_tree\_view\_controller} (\path{js/tree/file_tree_view_controller.js})
é o dono único do par \enquote{qual view está ativa} + o \emph{listener} do
botão de alternância (exatamente um, anexado uma vez). Ele:

\begin{itemize}[leftmargin=1.6em]
  \item mantém o nome da view ativa (\lstinline|_activeView|) e os dados de
    hierarquia (\lstinline|_hierarchyData|);
  \item registra um \emph{renderer} por view (\lstinline|registerRenderer|); ao
    ativar a view, invoca o renderer correspondente;
  \item expõe \lstinline|showFileMode()|, \lstinline|showHierarchyMode()| e
    \lstinline|showStandardMode()|, além dos predicados
    \lstinline|isShowingFileMode/Hierarchy/Standard|;
  \item assina o \file{ProjectStore} (\lstinline|_wireProjectStore|): a view
    \enquote{standard} só entra no ciclo com um projeto aberto, então o botão
    habilita ao abrir e desabilita ao fechar (deixando \enquote{standard} para
    não ficar numa árvore de pastas vazia).
\end{itemize}

O conjunto de views pelas quais o botão rotaciona é dinâmico
(\lstinline|_cyclableViews|): \enquote{verilog} sempre presente, \enquote{hierarchy}
apenas com dados compilados, \enquote{standard} apenas com projeto aberto. O
botão fica habilitado quando há mais de uma view; um clique avança para a
próxima view do ciclo (\lstinline|_nextView|, com \emph{wrap}). Os renderers
registrados são delegadores que buscam o gerente correto em tempo de chamada:
\enquote{verilog} chama \lstinline|window.projectTreeManager.renderTree()|,
\enquote{hierarchy} chama \lstinline|renderHierarchicalTree| do
\lstinline|CompilationModule| mais recente e \enquote{standard} chama
\lstinline|window.standardTreeRenderer.render()|.

\subsection{Renderer da view \enquote{verilog}}

A renderização do \emph{verilog picker} fica no \file{RenderMixin}
(\path{js/project/project_tree_render.js}), misturado em
\lstinline|ProjectTreeManager.prototype| por \path{file_mode.js}. O
\lstinline|renderTree()| é um \emph{reconciler keyed-por-path}: nunca faz
\lstinline|innerHTML = ''|; compara \lstinline|this.verilogFiles| (chaveado por
caminho) contra as linhas existentes — linhas cujo caminho sumiu são removidas,
linhas ainda desejadas são atualizadas \emph{in-place}
(\lstinline|_updateFileItem|), e caminhos sem linha são criados
(\lstinline|_createFileItem|) e posicionados via \lstinline|insertBefore|. É
idempotente: o mesmo input duas vezes não causa mutações na segunda chamada.

As linhas software (\file{.cmm} em \path{Software/}) não recebem botão de
remoção (são geridas pelo processador). Arquivos importados que sumiram do disco
são exibidos num cartão de aviso (\lstinline|_renderMissingFilesNotice|).

\subsection{Renderer da view \enquote{standard}}

O \file{standard\_tree\_render} (\path{js/tree/standard_tree_render.js})
desenha a árvore de pastas enraizada em \lstinline|window.currentProjectPath|,
ao modo de um explorador de arquivos genérico. As leituras são \emph{lazy}: só
a raiz é listada de início; os filhos de uma pasta são buscados
(\lstinline|electronAPI.getFolderFiles|) na primeira expansão e mantidos no DOM
para reexpansão instantânea. Caminhos expandidos são lembrados num
\lstinline|Set| (\lstinline|_expanded|) para sobreviver a um re-render. A
ordenação é \enquote{pastas primeiro, depois alfabética} (\lstinline|sortEntries|).
Arquivos ignorados (\lstinline|isIgnored|): \emph{dotfiles}, o próprio
\path{.spf} e blobs de configuração legados
(\lstinline|projectOriented.json|, \lstinline|processorConfig.json|,
\lstinline|fileOriented.json|). O método \lstinline|revealFolder| expande os
ancestrais até uma pasta-alvo e a destaca; \lstinline|expandAll| /
\lstinline|collapseAll| operam sobre toda a árvore.

\begin{nota}
Há um único caminho de abertura de arquivo, idêntico ao da view \enquote{verilog}
(\lstinline|readFile| $\rightarrow$ painel dividido em foco ou \emph{preview} do
\file{TabManager}; o duplo-clique promove a aba a permanente). Não há um segundo
\emph{handler} de abertura — a classe de bug que motivou a remoção da antiga
view \enquote{standard} não retorna.
\end{nota}

\section{Marcação de top-level e testbench}

A categoria \emph{synth} vs \emph{testbench} de cada arquivo Verilog é
\emph{auto-detectada} do conteúdo pelo classificador
(\path{js/project/verilog_classifier.ts}), não por uma marca manual. O
classificador é um \emph{scorer} com limiar \lstinline|TB_THRESHOLD = 3|: cada
sinal definitivo (dump de VCD via \lstinline|$dump*|, terminação de simulação
\lstinline|$finish|/\lstinline|$stop|, módulo sem portas) soma \lstinline|+3| e
já passa o limiar sozinho; o bloco \lstinline|initial| soma \lstinline|+2|
(forte, mas não classifica isolado, pois código sintetizável pode tê-lo); e os
sinais fracos (\lstinline|$display| e afins, atrasos procedurais como
\lstinline|#10|, hint de nome \lstinline|tb|/\lstinline|test|/\lstinline|testbench|)
somam \lstinline|+1| ou \lstinline|+2| cada. Em dúvida, o default seguro é
\emph{synthesizable}. Arquivos \file{.py} (cocotb) são sempre testbench
(decisão em \lstinline|_classifyAll| de \path{file_mode.js}, antes do
\emph{scorer}).

A escolha de \emph{top-level} e \emph{testbench top}, porém, é explícita do
usuário (menu de contexto na linha do arquivo, em \path{project_tree_actions.js}).
A marca \lstinline|isTopLevel| é relativa à categoria atual do arquivo: um synth
com \lstinline|isTopLevel=true| é o \emph{synth top}; um testbench com
\lstinline|isTopLevel=true| é o \emph{testbench top}. As marcações são
persistidas no \path{.spf} via \lstinline|SpfStore.update| (campos
\lstinline|topLevelFile| / \lstinline|testbenchFile| e as listas
\lstinline|synthesizableFiles| / \lstinline|testbenchFiles|), e a remoção de um
arquivo limpa essas referências quando apontavam para o caminho removido.

\section{Tela de boas-vindas e projetos recentes}
\label{sec:recentes}

Há \emph{duas} listas de recentes, com propósitos distintos:

\begin{description}[leftmargin=2em,style=nextline]
  \item[Tela de boas-vindas (renderer)] gerida por
    \file{RecentProjectsManager} (\path{js/project/recent_projects.js}),
    persistida em \lstinline|localStorage| sob a chave
    \lstinline|aurora-recent-projects| (até 10 entradas, cada uma com
    \lstinline|name|, \lstinline|path|, \lstinline|lastOpened|). Alimenta o
    componente \lstinline|<aurora-welcome>| (Lit), que renderiza a etapa
    Start/Recent, anima as linhas e mostra um \emph{preview} dos processadores
    de cada projeto (lido sob demanda de cada \path{.spf} por
    \lstinline|_readProcessors|, cacheado em \lstinline|p._procs|).
  \item[\emph{Jumplist} do Windows (processo principal)] gerida por
    \path{main/recents.js}, persistida num único JSON em \lstinline|userData|
    (\path{aurora-recents.json}), também limitada a 10 (\lstinline|MAX_RECENTS|).
    A API é \lstinline|read|, \lstinline|push| (insere no topo,
    deduplicado/normalizado), \lstinline|prune| (remove entradas cujo \path{.spf}
    não existe mais) e a constante \lstinline|MAX_RECENTS|.
\end{description}

A separação existe porque o processo principal não alcança o
\lstinline|localStorage|, e a lista no main permite limpar os resíduos
\enquote{Electron} de execuções antigas da \emph{jumplist} gerida pela shell do
Windows. Os recentes do main são atualizados em toda abertura/criação de
projeto, junto da reconstrução da \emph{jumplist} (\lstinline|rebuildJumpList|).
O handler \lstinline|list-recent-projects| devolve a lista do main após
\lstinline|prune|.

A \file{RecentProjectsManager} valida a existência do \path{.spf} antes de abrir
(\lstinline|checkProjectExists|); se sumiu, remove da lista e avisa o usuário.
Ao clicar numa linha, fecha as abas abertas e delega para o callback de abertura
injetado (\lstinline|loadProject|).

\section{Backups (zip sob demanda)}

A \tool{AURORA} produz backups \file{.zip} sob demanda (botão
\lstinline|backupFolderBtn| na toolbar, habilitado com projeto aberto). O
handler \lstinline|create-backup| (\path{main/ipc/files.js}):

\begin{enumerate}[leftmargin=1.8em]
  \item valida o \lstinline|folderPath| (\lstinline|safePath|);
  \item garante a pasta \path{Backup/} dentro do projeto e uma pasta de
    \emph{staging} temporária \path{backup_<timestamp>/};
  \item copia todo o conteúdo do projeto para a \emph{staging} (exceto a própria
    \path{Backup/} e a \emph{staging});
  \item gera \path{Backup/<nome>_<timestamp>.zip} via
    \lstinline|Compress-Archive| do PowerShell (nativo do Windows, sem depender
    de 7-Zip no PATH);
  \item remove a pasta de \emph{staging} em qualquer caminho (sucesso ou
    falha), devolvendo \lstinline|{ success, message }|.
\end{enumerate}

O \emph{quoting} para o PowerShell escapa aspas simples (\lstinline|'| para
\lstinline|''|), e o comando roda com \lstinline|-NoProfile -NonInteractive|.

\section{Estrutura de um projeto no disco}

O diagrama abaixo resume a estrutura típica de um projeto com um processador
\lstinline|cpu| e a pasta de backups:

\begin{lstlisting}
demo/                          <- pasta do projeto (== nome do .spf)
|-- demo.spf                   <- arquivo de projeto (JSON: metadata + structure)
|-- cpu/                       <- processador SAPHO
|   |-- Software/
|   |   |-- cpu.cmm            <- fonte C+-  (#PRNAME, #NUBITS, ...)
|   |   `-- cpu.asm            <- gerado (nao aparece na arvore)
|   |-- Hardware/
|   |   `-- cpu.v              <- Verilog sintetizavel (synth top)
|   `-- Simulation/
|       `-- cpu_tb.v           <- testbench (testbench top)
|-- (arquivos .v importados)   <- grupo "Imported" na arvore
`-- Backup/
    `-- demo_2026-06-21.zip    <- backup sob demanda
\end{lstlisting}

\begin{nota}
A categoria de cada \file{.v} (synth ou testbench) é decidida pelo conteúdo, não
pela subpasta: \lstinline|_discoverProcessorFiles| em \path{file_mode.js} varre
as três subpastas (\path{Hardware}, \path{Software}, \path{Simulation}) e aceita
\file{.v}/\file{.sv}/\file{.vh} como sintetizáveis e \file{.cmm} como software,
independentemente da pasta; o \lstinline|_classifyAll| posteriormente reclassifica
cada Verilog. O \file{.asm} é gerado e re-descoberto a cada \emph{load}, ficando
fora da árvore para não gerar laço com o \emph{watcher}.
\end{nota}
